% SSSUP-FONTE-STORICA: fonti-locali/originali/2014-2020.zip
% Anno 2017, esercizio 2. Testo storico; figura TikZ ricostruita in PSTricks.
% Verificare editorialmente la trasposizione prima di considerarla definitiva.
\renewcommand{\Anno}{2017}
\renewcommand{\NumeroProblema}{2}

\renewcommand{\Testo}{
Nick ha un orologio con le lancette, con le ore regolarmente numerate da $1$ a $12$, a cui tiene molto.  Il suo amico Brigante, per scherzo, sostituisce numero e tipo degli ingranaggi in modo tale che ogni lancetta continui a muoversi con velocità angolare costante ma, ahinoi, sconosciuta.

Per capire la velocità delle due lancette, Nick vuole controllare il loro moto fra due rintocchi consecutivi
del campanile di Pisa e osserva che:
\begin{itemize}
\item al primo rintocco, le due lancette non sono sovrapposte;
\item al rintocco successivo, quindi dopo esattamente un'ora,
\begin{itemize}
\item la lancetta più corta ha fatto più di $2$ ma meno di $3$ giri in senso orario;
\item le due lancette si sono sovrapposte $4$ volte.
\end{itemize}
\end{itemize}

\begin{enumerate}
\item[(a)] È possibile che la lancetta più lunga sia ferma?
\item[(b)] Dopo un'ora quanti giri può aver fatto la lancetta lunga?
\item[(c)] Durante un minuto (di tempo vero scandito dal campanile di Pisa), la lancetta più veloce di
  quanto può essersi spostata?
\end{enumerate}
}

\renewcommand{\Risposta}{}
\renewcommand{\Aiutino}{}
\renewcommand{\SoluzioneUfficiale}{}
\renewcommand{\SoluzioneRedazionale}{}
\renewcommand{\SoluzioniAlternative}{}
\renewcommand{\Osservazioni}{}
\renewcommand{\Commenti}{}
